\chapter{System Design: Authenticated AES Pipeline}

\section{Design Overview}

The final system is designed as a reproducible zkVM workflow where encryption,
authentication, proof generation, and benchmark evidence are treated as one
auditable computation pipeline rather than separate unconnected steps. Every
reported value should be traceable to configuration, guest execution, receipt
verification, and a concrete benchmark artifact.

At a high level, the pipeline executes authenticated encryption and proof
generation inside one host/guest interaction:

\begin{itemize}
\item encrypt plaintext with AES-CTR;
\item compute an HMAC-SHA256 tag (truncated to 192 bits);
\item commit outputs to the guest journal;
\item generate and verify a receipt for the complete execution.
\end{itemize}

Both confidentiality and integrity logic run in guest code so outputs are bound
to one verifiable execution trace.

Chapter 7 explains why this design was chosen; this chapter explains how it is
structured and how data moves through the implementation.

\section{Design Alternatives and Trade-offs}

Three designs were considered before the final split:

\begin{itemize}
\item \textbf{Host-only encryption:} cheapest to implement, but it would not prove
that ciphertext and tag were computed by the guest image.
\item \textbf{One guest target with runtime mode flag:} fewer guest binaries,
but each benchmark result would need extra metadata to show whether it ran CTR
only or CTR+HMAC.
\item \textbf{Separate CTR and CTR+HMAC targets:} more repository duplication, but
clearer traceability and cleaner evaluation.
\end{itemize}

The third design was selected because the dissertation needs auditable evidence
more than minimal code structure. Separate targets make it obvious which guest
binary generated each artifact, and they avoid attributing authentication cost to
the confidentiality-only path.

Figure~\ref{fig:host-guest-prover-verifier-flow} summarizes the selected host,
guest, prover, and verifier responsibilities.

\begin{figure}[t]
\centering
\fbox{%
\begin{minipage}{0.94\linewidth}
\small
\textbf{Host}
(serialize keys/IV/plaintext/AAD, choose benchmark target)
$\rightarrow$
\textbf{Guest}
(AES-CTR encryption, HMAC tag, journal commit)
$\rightarrow$
\textbf{Prover}
(receipt generation over execution trace)
$\rightarrow$
\textbf{Verifier}
(\texttt{METHOD\_ID} + journal decoding checks).
\end{minipage}}
\caption{Host $\rightarrow$ Guest $\rightarrow$ Prover $\rightarrow$ Verifier flow used by the authenticated pipeline.}
\label{fig:host-guest-prover-verifier-flow}
\end{figure}

\section{Encrypt-then-MAC Construction}

The implemented authenticated mode follows Encrypt-then-MAC:

\begin{itemize}
\item encryption call: \texttt{encrypt\_ctr(...)};
\item tag call: \texttt{compute\_auth\_tag\_192(...)};
\item verification path: \texttt{verify\_then\_decrypt(...)} checks tag before
CTR decryption.
\end{itemize}

This verify-before-decrypt order prevents processing unauthenticated ciphertext
inside the intended API path.

The ordering property matters for both security reasoning and implementation
hygiene. If decryption occurred before tag verification, malformed ciphertext
could influence downstream logic even when integrity ultimately fails. By
checking the tag first, the implementation keeps failure handling local and
prevents accidental dependence on unauthenticated plaintext bytes.

Key-separation assumptions are preserved by using distinct key material for
encryption and authentication, and by treating tag generation as a separate
primitive call. This is consistent with the composition model used in the
security analysis chapter.

Tag truncation to 192 bits is an explicit engineering choice in this project.
Chapter 10 includes the HMAC tag as part of the proof-output IND-CPA security
game.
In the standard ideal-PRF intuition, a single blind forgery attempt succeeds with
probability about \(2^{-192}\), and even many practical attempts remain far below
the risk scale considered in this dissertation. The selected tag length is
therefore shorter than full SHA-256 output while still leaving a very large
security margin for the scoped benchmark system.

\section{Execution Flow Across Host and Guest}

At runtime, the host serializes mode-specific inputs (keys, IV/nonce material,
plaintext, and AAD where applicable); the guest performs deterministic
cryptographic computation and commits
public outputs; the prover emits a receipt; and the host verifies the receipt
against the expected image ID and journal content.

This flow ensures that ``ciphertext + tag'' claims are attached to a concrete,
verifiable execution record, not just to a host-side API return value.

The design intentionally keeps deterministic cryptographic state transitions in
guest code and orchestration responsibilities in host code. The trade-off is that
host-side mistakes can still affect experiment setup: the host chooses input
sizes, target commands, and expected image IDs. Those choices are captured in
benchmark configuration and run manifests so measurement evidence can be checked
alongside receipt-level correctness evidence.

\section{Execution Modes and Benchmark IDs}

The current implementation uses a structural target split for mode selection:

\begin{itemize}
\item \texttt{aes-ctr}: CTR-only encryption benchmarks;
\item \texttt{aes-ctr-hmac}: CTR+HMAC authenticated benchmarks.
\end{itemize}

This replaces the older ``single target with runtime mode switch'' style. In the
current repository, comparison between confidentiality-only and authenticated
paths is made by selecting different benchmark targets and IDs, not by toggling
a single host flag.

The benchmark harness therefore tracks both families explicitly, for example:

\begin{itemize}
\item \texttt{aes-ctr-1blk}, \texttt{aes-ctr-4blk}, \texttt{aes-ctr-16blk},
\texttt{aes-ctr-64blk};
\item \texttt{aes-ctr-hmac-1blk}, \texttt{aes-ctr-hmac-4blk},
\texttt{aes-ctr-hmac-16blk}, \texttt{aes-ctr-hmac-64blk}.
\end{itemize}

This explicit naming improves result traceability because each benchmark ID now
encodes both mode family and payload size.

Current benchmark artifacts provide matched CTR-only and CTR+HMAC coverage at 1,
4, 16, and 64 blocks. This removes mismatched-coverage ambiguity when reporting
authentication overhead in Chapter 11.

\section{Host vs Guest Responsibility Boundary}

The placement rule is simple: security-relevant deterministic transformations
belong in the guest, while experiment orchestration belongs in the host.
Encryption, tag computation, tag verification, and journal commitments are guest
responsibilities because those are the claims a verifier should be able to bind
to the receipt. Input generation, command-line parsing, benchmark repetition,
and artifact persistence remain host responsibilities because they are policy and
deployment concerns rather than the cryptographic statement being proved.

This boundary keeps the proof statement small enough to reason about while still
capturing the important security behavior. As argued in Chapter 7, moving HMAC to
the host would reduce guest work but weaken the receipt-binding claim. Sender
identity, if required, remains an outer transport guarantee rather than part of
the proved cryptographic relation.
